\section{Exact restriction descent toward the M\"obius corner}
\label{sec:tpc44-coefficient-descent}

Conditional expectation deletes Walsh coefficients.  Physical
de--randomization does something different: setting a sign to \(-1\)
subtracts paired coefficients.  This section records that operation
exactly.

Let \(p_1<\cdots<p_R\) be as in \eqref{eq:tpc44-active-primes} and put
\begin{equation}
 P_r:=\prod_{j=1}^{r}p_j,
 \qquad P_0:=1.
 \label{eq:tpc44-Pr}
\end{equation}
For squarefree \(m\) with \((m,P_r)=1\), define
\begin{equation}
 C_r(m):=\sum_{a\mid P_r}\mu(a)C(am).
 \label{eq:tpc44-Cr-definition}
\end{equation}
Recall that \(C\) is extended by zero and \(C(1)=0\).
Throughout this section, every sum over \(m\) is restricted to squarefree
positive integers.  We also use the convention \(p_0=1\), so restriction
to coordinates \(p\le p_0\) means no restriction.

\begin{proposition}[Exact physical restriction]
\label{prop:tpc44-physical-restriction}
For every \(0\le r\le R\),
\begin{equation}
 F(-\one_{\le p_r},\varepsilon_{>p_r})
 =\sum_{\substack{m\ge1\\(m,P_r)=1}}C_r(m)\chi_m(\varepsilon).
 \label{eq:tpc44-restriction-expansion}
\end{equation}
The coefficients satisfy
\begin{equation}
 C_r(m)=C_{r-1}(m)-C_{r-1}(p_rm)
 \qquad (m,P_r)=1.
 \label{eq:tpc44-coefficient-recurrence}
\end{equation}
The physical signed sum is invariant under the descent:
\begin{equation}
 F(-\one)
 =\sum_{(m,P_r)=1}\mu(m)C_r(m),
 \label{eq:tpc44-physical-invariance}
\end{equation}
and at the terminal stage
\begin{equation}
 F(-\one)=C_R(1).
 \label{eq:tpc44-terminal-coefficient}
\end{equation}
\end{proposition}

\begin{proof}
Write every squarefree kernel uniquely as \(\kappa=am\), where
\(a\mid P_r\) and \((m,P_r)=1\).  At the restricted face,
\(
 \chi_a(-\one)=\mu(a)
\), so grouping by \(m\) proves \eqref{eq:tpc44-restriction-expansion}.
Splitting the divisors of \(P_r=P_{r-1}p_r\) according as they contain
\(p_r\) gives \eqref{eq:tpc44-coefficient-recurrence}.  Evaluating the
remaining signs at \(-1\) gives \eqref{eq:tpc44-physical-invariance}.
After every active prime has been processed, the only possible remaining
index is \(m=1\), proving \eqref{eq:tpc44-terminal-coefficient}.
\end{proof}

Define the restriction energy
\begin{equation}
 E_r:=\sum_{(m,P_r)=1}|C_r(m)|^2.
 \label{eq:tpc44-Er}
\end{equation}
Thus \(E_0=V\) and \(E_R=|F(-\one)|^2\).

\begin{theorem}[Exact coefficient--correlation product]
\label{thm:tpc44-correlation-product}
Whenever \(E_{r-1}>0\), define
\begin{equation}
 \theta_r:=
 \frac{
  2\operatorname{Re}
  \displaystyle\sum_{(m,P_r)=1}
  C_{r-1}(m)\overline{C_{r-1}(p_rm)}
 }{E_{r-1}}.
 \label{eq:tpc44-theta}
\end{equation}
Then
\begin{equation}
 -1\le\theta_r\le1,
 \qquad
 E_r=(1-\theta_r)E_{r-1}.
 \label{eq:tpc44-Er-recurrence}
\end{equation}
If an energy first becomes zero, set every later \(\theta_r=0\); every
later energy is then zero.  With this convention,
\begin{equation}
 \boxed{
 |F(-\one)|^2
 =V\prod_{r=1}^{R}(1-\theta_r).}
 \label{eq:tpc44-exact-correlation-product}
\end{equation}
\end{theorem}

\begin{proof}
Partition the indices in \(E_{r-1}\) into pairs \(m,p_rm\), with
\((m,P_r)=1\).  Then
\begin{align}
 E_{r-1}
 &=\sum_{(m,P_r)=1}
 \left(
  |C_{r-1}(m)|^2+|C_{r-1}(p_rm)|^2
 \right),
 \label{eq:tpc44-paired-old-energy}\\
 E_r
 &=\sum_{(m,P_r)=1}
 |C_{r-1}(m)-C_{r-1}(p_rm)|^2.
 \label{eq:tpc44-paired-new-energy}
\end{align}
Expanding the second line gives \eqref{eq:tpc44-Er-recurrence}.  The
inequality
\(
 2|ab|\le|a|^2+|b|^2
\)
gives \(|\theta_r|\le1\).  Iteration from \(E_0=V\) to
\(E_R=|F(-\one)|^2\) proves the product formula.
\end{proof}

The order of the prime restrictions can be changed.  The intermediate
\(\theta_r\) then change, but their product is forced by the endpoints.
Thus \eqref{eq:tpc44-exact-correlation-product} is an identity along a
chosen coordinate path, not an assertion that the local correlations
are independent.

We next relate the small--prime projection count to the restriction
energy.  Let \(\mathcal K_z\) be the set of low parts
\(\kappa_{\le z}\) of nonzero pair kernels.  Since every pair kernel is
the symmetric difference of two row labels,
\begin{equation}
 |\mathcal K_z|\le\nu_z^2.
 \label{eq:tpc44-kernel-projection-count}
\end{equation}

\begin{proposition}[Energy cost of freezing the low primes]
\label{prop:tpc44-low-restriction-energy}
After every active prime \(p\le z\) has been restricted to \(-1\), the
remaining coefficient energy \(E_z=E_{r(z)}\) satisfies
\begin{equation}
 E_z\le\nu_z^2V.
 \label{eq:tpc44-low-energy-bound}
\end{equation}
\end{proposition}

\begin{proof}
For a fixed high kernel \(m\), the restricted coefficient is a signed
sum
\begin{equation}
 C_{\le z}(m)
 =\sum_{a\in\mathcal K_z}\mu(a)C(am),
 \label{eq:tpc44-low-kernel-sum}
\end{equation}
with absent coefficients interpreted as zero.  Cauchy--Schwarz gives
\begin{equation}
 |C_{\le z}(m)|^2
 \le|\mathcal K_z|\sum_{a\in\mathcal K_z}|C(am)|^2.
 \label{eq:tpc44-low-kernel-Cauchy}
\end{equation}
Summing over \(m\), using unique low--high factorization, and applying
\eqref{eq:tpc44-kernel-projection-count} proves the claim.
\end{proof}

Although the exponent \(\nu_z^2\) is natural for the squared coefficient
spectrum, \cref{thm:tpc44-robust-face} is stronger for the nonnegative
energy itself: it returns to the unsquared row labels and pays only one
factor \(\nu_z\).
